\chapter{Criterio de anexo metodológico de claims observados}

\section{Anexo metodológico}
\begin{definition}
Se llama anexo metodológico de claims observados al espacio disciplinado donde se conservan claims que ya han sido formulados y revisados en alguna medida, pero que todavía requieren reserva, observación crítica o reformulación antes de cualquier salida pública fuerte.
\end{definition}

\begin{observation}
El anexo metodológico no funciona como descarte ni como depósito indiferenciado. Su función es preservar trazabilidad y valor metodológico de claims que aún no conviene proyectar como piezas externas estabilizadas.
\end{observation}

\section{Condición mínima de entrada}
\begin{definition}
Un claim puede entrar en el anexo metodológico cuando concurre al menos una de las siguientes situaciones:
\begin{enumerate}
  \item presenta fragilidad crítica no resuelta;
  \item su alcance todavía requiere reformulación;
  \item la relación entre evidencia y formulación sigue siendo insuficientemente sobria;
  \item su valor principal es interno, comparativo o de reserva metodológica;
  \item conviene conservarlo sin elevarlo todavía a salida pública.
\end{enumerate}
\end{definition}

\begin{observation}
La entrada al anexo no implica inutilidad. Puede indicar, precisamente, que el claim tiene valor suficiente como para no perderse, pero aún no ha alcanzado la forma adecuada para su proyección externa.
\end{observation}

\section{Anexo y trazabilidad}
\begin{definition}
Todo claim observado que entre en el anexo debe conservar origen legible, dependencia canónica visible, estado editorial y razón explícita de observación o reserva.
\end{definition}

\begin{observation}
Esta cláusula evita que el anexo se vuelva opaco o que las reservas pierdan continuidad con la tabla madre del Tomo III.
\end{observation}

\section{Anexo y tabla madre}
\begin{definition}
El anexo metodológico debe poder leerse como prolongación disciplinada de la tabla madre de demos y claims, manteniendo vínculo visible entre el claim observado y su posición previa dentro de la cartografía curatorial.
\end{definition}

\begin{observation}
El claim observado no desaparece del mapa; cambia de estatuto editorial y metodológico.
\end{observation}

\section{Condición mínima de salida del anexo}
\begin{definition}
Un claim puede salir del anexo metodológico cuando la observación crítica haya sido resuelta o cuando la reformulación alcanzada permita devolverlo a la tabla principal o a una ruta sobria de salida pública.
\end{definition}

\begin{observation}
La salida del anexo no debe depender de deseo editorial aislado, sino de mejora disciplinada visible.
\end{observation}

\section{Cláusula de alcance}
\begin{observation}
El criterio de anexo metodológico de claims observados fija una reserva disciplinada dentro de Tomo III. No constituye todavía una teoría completa de descarte final ni una taxonomía exhaustiva de todas las reservas posibles del corpus.
\end{observation}

